% Tipografía de los cuadernillos de Nimbo. La obra elige el papel y la tinta.
\setlength{\parindent}{0pt}
\setlength{\emergencystretch}{2em}
\raggedbottom
\renewcommand{\contentsname}{Índice de fichas}
\newcommand{\portadaglosario}[2]{%
  \thispagestyle{empty}\vspace*{.16\textheight}
  {\small\sffamily\color{accent}GLOSARIO\par}\medskip
  {\LARGE\bfseries\raggedright #1\par}\bigskip
  {\color{accent}\rule{22mm}{1pt}}\par\vfill
  #2\par\clearpage}
\newcommand{\abrirficha}[2]{%
  \clearpage{\small\sffamily\color{accent}#1\par}\smallskip
  {\LARGE\bfseries\raggedright #2\par}
  \addcontentsline{toc}{section}{#2}%
  \smallskip{\color{accent}\rule{\linewidth}{.5pt}}\par\medskip}
\newcommand{\campoficha}[2]{%
  \par\addvspace{5pt}%
  {\small\sffamily\bfseries\color{accent}#1\par}\nopagebreak[4]%
  {\raggedright #2\par}}
\newcommand{\datoficha}[2]{%
  \par\addvspace{3pt}%
  {\raggedright{\small\sffamily\bfseries\color{accent}#1}\enspace #2\par}}
\newcommand{\retratoficha}[1]{%
  \par{\centering\includegraphics[width=.8\linewidth,height=.29\textheight,keepaspectratio]{#1}\par}\smallskip}
